% Updated prospective human study (GCS-verified complete 6-baseline cohort).
% Distinct from frozen historical tex/humanstudy.tex; no private rows embedded.
% Build current trajectory/endpoint CSVs with:
% python3 human_study/scripts/render_humanstudy_new_pgfplots.py \
%   --analysis-dir ~/private/dtag-analysis-2026-10-30-restored
\subsection{Protocol and outcome}
We recruited participants through Prolific for the prospective DTAG
human-response experiment. Each calibrated session comprised six
baseline questions, twenty adaptive questions, and five interleaved
attention checks. Participants answered in natural language; accepted
native survey-variable assignments updated the LSM state. Four
question-selection policies were used: Left, Right, Neutral, and Random.
The original study randomized policies; subsequent cohorts used fixed
Left or Right assignments. Pooling across recruitment waves is
descriptive, not a randomized comparison between waves.

Let $s_0$ be the state after the six baseline questions and $s_t$
the state following adaptive answer $t$. The index trajectory and
baseline-relative displacement are
\begin{equation}
\Delta I_\Phi(t)=I_\Phi(s_t)-I_\Phi(s_0),
\qquad t=0,\ldots,20.
\label{eq:humanstudy-new-delta}
\end{equation}
A directional success at the end of the adaptive sequence is $\Delta I_\Phi(20)<0$ for Left
and $\Delta I_\Phi(20)>0$ for Right. The index is the
LSM-inferred survey-response phenotype, not a measured change
in underlying political beliefs.

\subsection{Completed sessions and directional response}
After archival recovery, the calibrated cohort contains 71
complete 20-question sessions: 25 from the original four-policy
experiment, 15 from the fixed-Left follow-up, and 31 from the
fixed-Right follow-up. The assigned-policy totals are Left
($n=21$), Right ($n=38$), Neutral ($n=7$), and Random ($n=5$).
The separately recruited zero-calibration Right cohort is not
combined with these participants, since it does not have an
individually measured six-question baseline.

The archived results show leftward displacement in 20/21
Left-push participants (95.2\%) and rightward displacement in
21/38 Right-push participants (55.3\%). Baseline position
strongly distinguishes the observed Right-push response:
12/29 initially Left-leaning participants (41.4\%) moved right,
compared with 9/9 initially Right-leaning participants (100\%).
For Left push, 16/17 initially Left-leaning participants
(94.1\%) and 4/4 initially Right-leaning participants (100\%)
moved left.

\begin{table*}[!t]
\centering
\caption{Baseline-stratified directional outcomes after the full adaptive
sequence. The last column also includes \emph{correct early identification of resistance}
for initially Left-leaning participants assigned Right push; it is not
an intervention-only success rate.}
\label{tab:humanstudy-new-stratified}
\small
\setlength{\tabcolsep}{6pt}
\begin{tabular}{lrrr}
\toprule
Initial inferred phenotype & \shortstack{Left: direction\\achieved} &
\shortstack{Right: direction\\achieved} & \shortstack{Right: movement or\\correct resistance}\\
\midrule
Initially Left & 16/17 (94.1\%) & 12/29 (41.4\%) & 29/29 (100\%)\\
Initially Right & 4/4 (100\%) & 9/9 (100\%) & 9/9 (100\%)\\
\midrule
Overall & 20/21 (95.2\%) & 21/38 (55.3\%) & 38/38 (100\%)\\
\bottomrule
\end{tabular}
\end{table*}

% The endpoint means/CI were not supplied as aggregates. Regenerate
% these directly from the 71 archived complete trajectories rather than
% silently recycling the published 49-session curves.
\IfFileExists{Figures/humanstudy_new_table.tex}{%
  % Auto-generated from restored private GCS final archives; 71 calibrated 20/20.
\begin{table*}[!t]
\centering
\caption{Human-study results by assigned policy. Endpoint displacement is
measured from each participant's six-question baseline; intervals are
95\% participant-bootstrap percentile intervals (5,000 replicates).}
\label{tab:humanstudy-participants}
\footnotesize
\setlength{\tabcolsep}{6pt}
\begin{tabular}{lrrrrr}
\toprule
Policy & $n$ & Baseline $I_\Phi<0$ & Baseline $I_\Phi>0$ & Mean $\Delta I_\Phi(20)$ (95\% CI) & Directional success \\
\midrule
Left & 21 & 17 & 4 & $-0.0989\;[-0.1200,-0.0741]$ & 20/21 (95.2\%) \\
Right & 38 & 29 & 9 & $+0.0225\;[+0.0006,+0.0452]$ & 21/38 (55.3\%) \\
Neutral & 7 & 6 & 1 & $-0.0054\;[-0.0153,+0.0032]$ & --- \\
Random & 5 & 2 & 3 & $-0.0045\;[-0.0512,+0.0413]$ & --- \\
\bottomrule
\end{tabular}
\end{table*}
%
}{%
  \PackageError{DTAG}{Current human-study endpoint table not generated}
  {Run render_humanstudy_new_pgfplots.py on the restored private analysis export.}%
}

\begin{figure*}[!t]
\centering
\DTAGExternalEnable
% NEW 71-completer six-calibration cohort. Do not use old humanstudy CSVs here.
% Human DTAG empirical panels. All coordinates are read from regenerated 71-session calibrated aggregate CSVs.
% Generated data: python3 human_study/scripts/render_humanstudy_new_pgfplots.py
% Compile from within tex/ using dtag_main_new.tex.
% Styles follow DTAG figures 1--2: publication axis, bold labels and strong ticks.
%% {\small\sffamily\bfseries (a) Adaptive response trajectories}\par\vspace{1.0mm}
\begin{tikzpicture}[font=\bf\sffamily\fontsize{9}{9}\selectfont]
  \node[] (A) at (0,0) {
    \begin{tikzpicture}[anchor=center]
\begin{axis}[
  publication axis,
  width=2.2in, height=2in,
  scale only axis=true,
  xmin=0, xmax=20,
  %xtick={0,5,10,15,20},
  xlabel={Adaptive questions answered},
  ylabel={Change in inferred ideology index},
  axis line style={black,line width=1pt},
  tick style={black,line width=.95pt},
  tick label style={font=\footnotesize\bfseries},
  label style={font=\small\bfseries},
  title style={font=\small\bfseries},
  legend style={draw=none,fill=none,at={(0.0,.77)},anchor=south west},
  legend cell align=left, legend columns=2,
  ymajorgrids, grid style={black!12,line width=.3pt},title={(a) Adaptive response trajectories},
]
\addplot[black!45,densely dashed,thin,forget plot] coordinates {(0,0) (20,0)};
% Participant-level 95% confidence bands only for Left and Right.
\addplot[name path=hsleftlo,draw=none,forget plot]
  table[x=step,y=ci_low,col sep=comma]{Figures/humanstudy_new_trajectory_left.csv};
\addplot[name path=hslefthi,draw=none,forget plot]
  table[x=step,y=ci_high,col sep=comma]{Figures/humanstudy_new_trajectory_left.csv};
\addplot[fill=dtblue,fill opacity=.14,draw=none,forget plot]
  fill between[of=hsleftlo and hslefthi];
\addplot[name path=hsrightlo,draw=none,forget plot]
  table[x=step,y=ci_low,col sep=comma]{Figures/humanstudy_new_trajectory_right.csv};
\addplot[name path=hsrighthi,draw=none,forget plot]
  table[x=step,y=ci_high,col sep=comma]{Figures/humanstudy_new_trajectory_right.csv};
\addplot[fill=dtred,fill opacity=.14,draw=none,forget plot]
  fill between[of=hsrightlo and hsrighthi];
\addplot[dtblue,ultra thick]
  table[x=step,y=mean,col sep=comma]{Figures/humanstudy_new_trajectory_left.csv};
\addlegendentry{Left}
\addplot[dtred,ultra thick]
  table[x=step,y=mean,col sep=comma]{Figures/humanstudy_new_trajectory_right.csv};
\addlegendentry{Right}
\addplot[black!75,ultra thick,densely dashed]
  table[x=step,y=mean,col sep=comma]{Figures/humanstudy_new_trajectory_neutral.csv};
\addlegendentry{Neutral}
\addplot[modteal,ultra thick]
  table[x=step,y=mean,col sep=comma]{Figures/humanstudy_new_trajectory_random.csv};
\addlegendentry{Random}
\end{axis}
\end{tikzpicture}};
%% \end{minipage} 
%% \hfill
%% \begin{minipage}[t]{0.365\textwidth}
%% \centering
%% {\small\sffamily\bfseries (b) Twenty-question endpoints}\par\vspace{1.0mm}
\node[anchor=south west] (B) at (A.south east) {\begin{tikzpicture}[anchor=center]
\begin{axis}[
  publication axis,
  width=2in, height=2.2in,
  scale only axis=true,
  xmin=.45, xmax=4.55,
  xtick={1,2,3,4},
  xticklabels={Left,Right,Neutral,Random},
  x tick label style={rotate=0},
  ylabel={Mean endpoint $\Delta$ (95\% CI)},
  axis line style={black,line width=1pt},
  tick style={black,line width=.95pt},
  tick label style={font=\footnotesize\bfseries},
  label style={font=\small\bfseries},
  title style={font=\small\bfseries},
  label style={},
  scaled ticks=false,
  ymajorgrids, grid style={black!12,line width=.3pt},title={(b) Twenty-question endpoints}
]
\addplot[black!45,densely dashed,thin,forget plot] coordinates {(.5,0) (4.5,0)};
% Mean +/- participant-bootstrap 95% CI for each of all four arms.
\addplot+[only marks,mark=*,mark size=2.5pt,color=dtblue,
  error bars/.cd,y dir=both,y explicit,error bar style={line width=1.05pt},
  error mark options={mark size=2pt,line width=1.05pt, rotate=90}]
  table[x=x,y=mean,y error minus=err_minus,y error plus=err_plus,col sep=comma]
    {Figures/humanstudy_new_endpoint_left.csv};
\addplot+[only marks,mark=*,mark size=2.5pt,color=dtred,
  error bars/.cd,y dir=both,y explicit,error bar style={line width=1.05pt},
  error mark options={mark size=2pt,line width=1.05pt, rotate=90}]
  table[x=x,y=mean,y error minus=err_minus,y error plus=err_plus,col sep=comma]
    {Figures/humanstudy_new_endpoint_right.csv};
\addplot+[only marks,mark=*,mark size=2.5pt,color=black!75,
  error bars/.cd,y dir=both,y explicit,error bar style={line width=1.05pt},
  error mark options={mark size=2pt,line width=1.05pt, rotate=90}]
  table[x=x,y=mean,y error minus=err_minus,y error plus=err_plus,col sep=comma]
    {Figures/humanstudy_new_endpoint_neutral.csv};
\addplot+[only marks,mark=*,mark size=2.5pt,color=modteal,
  error bars/.cd,y dir=both,y explicit,error bar style={line width=1.05pt},
  error mark options={mark size=2pt,line width=1.05pt, rotate=90}]
  table[x=x,y=mean,y error minus=err_minus,y error plus=err_plus,col sep=comma]
    {Figures/humanstudy_new_endpoint_random.csv};
\end{axis}
\end{tikzpicture}};
\end{tikzpicture}


\vspace{-8pt}
\caption{\textbf{Updated human-study trajectories and endpoints
(71 calibrated complete sessions).}
(a) Mean baseline-relative change after each adaptive question in
Left, Right, Neutral, and Random policies; the colored 95\%
participant-bootstrap confidence bands are shown for Left and Right.
(b) Mean displacement after the twenty adaptive questions and 95\% participant-bootstrap intervals,
computed from complete archived sessions. All trajectories in this
figure use only GCS-verified 20-answer completers, rather than
imputing values for participants who stopped early.}
\label{fig:humanstudy-new-trajectory-endpoints}
\end{figure*}

\subsection{Early identification of Right-push resistance}
Right-directed questioning produced a weaker response among initially
Left-leaning participants. We therefore examined whether resistance
to rightward displacement could be identified from the first four
adaptive responses. This analysis was restricted to the 29 Right-push
participants whose baseline inferred ideology index was negative.
A participant was flagged as resistant if their inferred index had
not increased relative to their own baseline after four questions.
We evaluated that early classification against whether the index
had moved rightward from baseline at the end of the full
20-question sequence.

Seventeen of these 29 participants showed no net rightward
displacement at the end of the study. All 17 had been identified
by the four-question criterion. The criterion also flagged two
participants who ultimately moved rightward. It therefore
identified all observed resistant participants (17/17,
100\% sensitivity), with 89.5\% precision (17/19),
83.3\% specificity (10/12), and 93.1\% overall
classification accuracy (27/29).

\begin{table}[!t]
\centering
\caption{Early resistance classification among initially Left-leaning
participants assigned Right push ($n=29$). The early classification
uses the first four adaptive questions; the outcome is net movement
relative to baseline at the end of the 20-question study.}
\label{tab:humanstudy-q4-confusion}
\footnotesize
\setlength{\tabcolsep}{3pt}
\begin{tabular}{lcc}
\toprule
 & \shortstack{No net\\rightward move} & \shortstack{Moved\\rightward}\\
\midrule
Flagged resistant & 17 & 2 \\
Not flagged & 0 & 10 \\
\bottomrule
\end{tabular}
\end{table}

\begin{figure*}[!t]
\centering
% Updated DTAG human-study response/resistance panels, 2026-10-30.
% Aggregate inputs only; no participant trajectories or identifiers.
% Both panels preserve the manuscript's thick axes, bold labels and ticks,
% dtblue/dtred palette and restrained grid lines.
\begin{tikzpicture}[font=\bf\sffamily\fontsize{9}{9}\selectfont]
\node (A) at (0,0) {
\begin{tikzpicture}
\begin{axis}[
  publication axis,
  width=2.9in,height=2.15in,scale only axis=true,
  xmin=.5,xmax=4.5,ymin=0,ymax=115,
  xtick={1,2,3,4},
  xticklabels={L/L,L/R,R/L,R/R},
  tick label style={font=\footnotesize\bfseries},
  x tick label style={font=\footnotesize\bfseries},
  label style={font=\small\bfseries},
  title style={font=\small\bfseries},
  ylabel={Directional response (\%)},
  xlabel={Push / initial phenotype},
  ytick={0,25,50,75,100},
  title={(c) Directional movement by baseline},
  axis line style={black,line width=1pt},
  tick style={black,line width=.95pt},
  ymajorgrids,grid style={black!12,line width=.3pt}
]
\addplot+[ybar,mark=none,bar width=15pt,draw=dtblue,fill=dtblue!80,
  nodes near coords={\pgfmathprintnumber[fixed,precision=1]{\pgfplotspointmeta}\%},
  point meta=y,every node near coord/.append style={font=\scriptsize\bfseries},
  line width=.9pt]
 table[x=x,y=percent,col sep=comma]{Figures/humanstudy_new_success_left.csv};
\addplot+[ybar,mark=none,bar width=15pt,draw=dtred,fill=dtred!80,
  nodes near coords={\pgfmathprintnumber[fixed,precision=1]{\pgfplotspointmeta}\%},
  point meta=y,every node near coord/.append style={font=\scriptsize\bfseries}]
 table[x=x,y=percent,col sep=comma]{Figures/humanstudy_new_success_right.csv};
\end{axis}
\end{tikzpicture}};
\node[anchor=south west] (B) at ([xshift=8mm]A.south east) {
\begin{tikzpicture}
\begin{axis}[
  publication axis,
  width=2.85in,height=2.15in,scale only axis=true,
  xmin=.5,xmax=4.5,ymin=0,ymax=20,
  xtick={1,2,3,4},
  xticklabels={TP,FP,FN,TN},
  tick label style={font=\footnotesize\bfseries},
  label style={font=\small\bfseries},
  title style={font=\small\bfseries},
  ylabel={Participants},
  xlabel={Q4 prediction versus Q20 outcome},
  ytick={0,5,10,15,20},
  title={(d) Q4 Right-resistance classification},
  axis line style={black,line width=1pt},
  tick style={black,line width=.95pt},
  ymajorgrids,grid style={black!12,line width=.3pt}
]
\addplot+[ybar,mark=none,bar width=15pt,draw=modteal,fill=modteal!75,
  nodes near coords,point meta=y,
  every node near coord/.append style={font=\footnotesize\bfseries}]
 table[x=x,y=n,col sep=comma]{Figures/humanstudy_new_q4_correct.csv};
\addplot+[ybar,mark=none,bar width=15pt,draw=dtred,fill=dtred!75,
  nodes near coords,point meta=y,
  every node near coord/.append style={font=\footnotesize\bfseries}]
 table[x=x,y=n,col sep=comma]{Figures/humanstudy_new_q4_errors.csv};
\end{axis}
\end{tikzpicture}};
\end{tikzpicture}

\vspace{-8pt}
\caption{\textbf{Direction-dependent response and early resistance detection.}
(c) Proportion moving in the assigned direction after the study by
initial phenotype sign and policy: L/L $=16/17$, L/R $=4/4$,
R/L $=12/29$, and R/R $=9/9$, where the first symbol denotes
policy and the second denotes baseline sign.
(d) Classification from the first four questions of baseline-relative
Right-push nonresponse at the end of the study among initially Left-leaning participants:
17 true positives (TP), 2 false positives (FP), no false negatives
(FN), and 10 true negatives (TN).}
\label{fig:humanstudy-new-response-resistance}
\end{figure*}

Rightward displacement or correctly identified endpoint
nonresponse accounts for all 29 initially Left-leaning Right-push
participants (12 and 17, respectively); including initially
Right-leaning participants, the combined outcome is 38/38.
This is a composite of successful directional response and
correct early identification of nonresponse, \emph{not} a
38/38 directional intervention success rate.

\subsection{Interpretation}
The completed human study shows asymmetric baseline-relative
movement in the inferred index. Left-directed question selection
produced the intended signed change in 20/21 sessions. Right
push was effective for all nine initially Right-leaning
participants but for only 12/29 initially Left-leaning
participants. In the latter group, all observed endpoint
nonresponders were identified using the first four adaptive
responses. These findings establish the descriptive relationship
between initial inferred position, directional response,
and an early baseline-relative resistance flag in this
experiment; they do not establish persuasion or changes in
participants' underlying political beliefs.
